\chapter{Flutter of a continuous wing with nothing discretized}\label{ch:wing}

Chapter~\ref{ch:aero} treated a wing \emph{section}: two degrees of freedom
and exact Theodorsen aerodynamics. A wing is a flexible beam with infinitely
many modes. The standard route to its flutter involves two independent
approximations:
\begin{enumerate}
  \item a spatial discretization, usually finite elements, followed by a
        truncation to $N$ vibration modes;
  \item a rational fit of the unsteady aerodynamic forces, tabulated at a
        few reduced frequencies, which turns the system into a large
        state-space model \citep{edwards1977,kaiser2022}.
\end{enumerate}
Each approximation raises its own convergence question, and the fit can
introduce spurious poles. This chapter computes the poles, the stability,
the flutter speed and the time responses of the classical Goland wing
\citep{goland1945} with \emph{neither} approximation. The bending and
torsion equations are solved exactly along the span, and Theodorsen's
function is evaluated from Bessel functions at every complex $s$. The only
modelling assumption on the aerodynamics is strip theory. The code is in
\path{examples/continuous_wing/} and \path{examples/models/wing_*.m}.

\section{Model}

\begin{figure}[H]
\centering
\begin{tikzpicture}[>=Latex,thick,font=\small,x=1cm,y=1cm]
  % Planform: span along x (0 = root, 9 = tip), chord along y (LE at top).
  \fill[blue!8] (0,-1) rectangle (9,1);
  \draw[blue!60!black] (0,-1) rectangle (9,1);
  % Clamped root
  \fill[pattern=north east lines] (-0.35,-1.5) rectangle (0,1.5);
  \draw (0,-1.5)--(0,1.5);
  % Elastic axis (33% chord) and centre of gravity (43% chord)
  \draw[orange!80!black,dashed] (0,0.34)--(9,0.34);
  \node[orange!80!black,font=\scriptsize,anchor=west] at (9.1,0.34) {elastic axis};
  \draw[black,dotted] (0,0.14)--(9,0.14);
  \node[font=\scriptsize,anchor=west] at (9.1,0.0) {centre of mass};
  \node[font=\scriptsize,anchor=west] at (9.1,1.0) {leading edge};
  % One strip with exact Theodorsen loads
  \fill[teal!25] (5.4,-1) rectangle (6.1,1);
  \draw[teal!60!black] (5.4,-1) rectangle (6.1,1);
  \draw[->,teal!60!black,very thick] (5.75,0.55)--(5.75,1.75);
  \node[teal!50!black,font=\scriptsize,align=center] at (5.75,2.35)
      {strip loads $Q(s,U)\,[w;\alpha]$\\exact $C(sb/U)$, no fit};
  % Free stream
  \foreach \x in {1.2,2.6,4.0} \draw[->,blue!60!black] (\x,2.1)--(\x,1.25);
  \node[blue!60!black] at (2.6,2.35) {$U$};
  % Degrees of freedom at the tip (section view arrows)
  \draw[->,red!70!black] (8.55,0.34)--(8.55,-0.6);
  \node[red!70!black,font=\scriptsize,anchor=east] at (8.5,-0.45) {$w$ (down)};
  \draw[->,red!70!black] (7.3,0.75) arc[start angle=150,end angle=30,radius=0.45];
  \node[red!70!black,font=\scriptsize,anchor=south] at (7.7,0.95) {$\alpha$ (nose-up)};
  % Span coordinate
  \draw[->] (0,-1.9)--(9.4,-1.9) node[right] {$y$};
  \draw (9,-1.8)--(9,-2.0) node[below] {$L$};
  \node[below] at (0,-2.0) {$0$};
  % Boundary conditions and the exact propagation
  \node[font=\scriptsize,align=center] at (0.3,-2.75)
      {clamped root\\$w=w'=\alpha=0$};
  \node[font=\scriptsize,align=center] at (8.8,-2.75)
      {free tip\\$V=M=T=0$};
  \node[font=\scriptsize,align=center] at (4.5,-3.75)
      {$z=[w,\,w',\,\alpha,\,V,\,M,\,T]^T,\qquad z'=A(s,U)\,z
       \;\Rightarrow\; z(L)=e^{A(s,U)L}\,z(0)$\\[2pt]
       no finite elements, no modal truncation};
\end{tikzpicture}

\caption{The continuous wing. Each spanwise strip carries the exact
two-dimensional Theodorsen load; the state $z(y)$ is propagated exactly
from the clamped root to the free tip.}
\label{fig:wing}
\end{figure}

The wing (Figure~\ref{fig:wing}) is a straight cantilever of span $L$. Its
section at position $y$ moves with the plunge $w(y,t)$, positive
downwards at the elastic axis, and the nose-up pitch $\alpha(y,t)$. Bending
follows Euler--Bernoulli theory (stiffness $EI$) and torsion Saint-Venant
theory (stiffness $GJ$). The mass per unit span is $\mu$, the pitch
inertia $I_\alpha$ and the static moment $S=\mu x_\alpha b$; the static
moment couples the two motions through inertia. For a motion $e^{st}$, the
section equations are
\begin{equation}
  EI\,w''''+s^2(\mu w+S\alpha)=Q_{11}w+Q_{12}\alpha,\qquad
  -GJ\,\alpha''+s^2(Sw+I_\alpha\alpha)=Q_{21}w+Q_{22}\alpha,
  \label{eq:wingpde}
\end{equation}
with $'=d/dy$. $Q(s,U)$ is the strip load matrix of
Chapter~\ref{ch:aero}, $Q=\pi\rho U^2(p^2A_2+pA_1+A_0)$ with $p=sb/U$ and
Theodorsen's function $C(p)=K_1(p)/(K_0(p)+K_1(p))$
\citep{theodorsen1935,kaiser2022}. The example uses the Goland parameters
as tabulated by \citet[Table~2.4]{cal1992}:
\begin{itemize}
  \item $L=6.096$~m, $b=0.9144$~m;
  \item $EI=9.773\times10^6$~N\,m$^2$, $GJ=9.876\times10^5$~N\,m$^2$;
  \item $\mu=35.717$~kg/m, $I_\alpha=8.642$~kg\,m;
  \item elastic axis at 33\,\% of the chord, centre of mass at 43\,\%;
  \item $\rho=1.225$~kg/m$^3$.
\end{itemize}

\section{An exact characteristic function}

With $s$ as a parameter, \eqref{eq:wingpde} is a linear ordinary
differential equation in $y$. Collect displacement, slope, twist, shear
force, bending moment and torque in the state vector
\begin{equation}
  z=[\,w,\;w',\;\alpha,\;V,\;M,\;T\,]^T,\qquad M=EI\,w'',\quad V=-M',\quad T=GJ\,\alpha'.
\end{equation}
Then $z'=A(s,U)\,z$, where the only nonzero entries of $A$ are
\begin{equation}
  A_{12}=1,\; A_{25}=\tfrac1{EI},\; A_{36}=\tfrac1{GJ},\; A_{54}=-1,\qquad
  [A_{41}\;A_{43}]=D_{1,:},\; [A_{61}\;A_{63}]=D_{2,:},
\end{equation}
with $D(s,U)=s^2\bigl[\begin{smallmatrix}\mu&S\\S&I_\alpha\end{smallmatrix}\bigr]-Q(s,U)$.
On a span of constant properties, the solution is exact:
$z(L)=e^{A(s,U)L}z(0)$. For a wing whose properties vary, the span is
divided into $N$ strips of constant properties, with one propagator
$E_j=e^{A_j\ell_j}$ each.

Let $x=[z_0;\dots;z_N]$ be the (fixed-scaled) states at the strip
boundaries. The clamped root, the $N$ propagation relations and the free tip
give a square homogeneous system $K(s,U)\,x=0$:
\begin{equation}
  w_0=w_0'=\alpha_0=0,\qquad z_{j}-E_j\,z_{j-1}=0\;(j=1,\dots,N),\qquad
  V_N=M_N=T_N=0.
\end{equation}
The wing has a mode at $s$ exactly when $K(s,U)$ is singular, so the
characteristic function is
\begin{equation}
  \Delta(s,U)=\det K(s,U).
  \label{eq:wingdelta}
\end{equation}
Three properties make \eqref{eq:wingdelta} suitable for the argument
principle:
\begin{itemize}
  \item The matrix exponential is an entire function of the entries of
        $A$. $\Delta$ is therefore entire in vacuum, and analytic except on
        the negative real axis (the branch cut of $C$) with air. So
        \code{AssumeAnalytic} holds in any rectangle that avoids that axis.
  \item Nothing is inverted. Condensing $K$ to a small transfer or hybrid
        matrix would divide by functions of $s$ and create artificial
        poles.
  \item For identical strips $e^{AL}=(e^{AL/N})^N$, so $\Delta$ does not
        depend on $N$. The study confirms this to $1.4\times10^{-14}$ for
        $N=1,2,4,8$. For a uniform wing the strips are bookkeeping, not
        discretization.
\end{itemize}
The state is scaled by constants, independent of $s$, to improve the
conditioning of $K$.

\section{Stability, root locus and flutter}

In vacuum and without inertial coupling, the zeros of $\Delta$ found by
\code{croots} reproduce the closed-form cantilever frequencies
($\beta_n^2\sqrt{EI/\mu}/L^2$ in bending, $(2n-1)\pi\sqrt{GJ/I_\alpha}/(2L)$
in torsion). The eight frequencies below 900~rad/s agree to
$4\times10^{-14}$. A finite-element model reaches this accuracy only in the
limit of a fine mesh.

Stability is decided, as in Chapter~\ref{ch:aero}, by counting the zeros
of $\Delta$ in a rectangle of the right half-plane. This rectangle
includes the real axis, so a divergence root would also be counted. In
$[10^{-3},60]\times[-400,400]$ the count is 0 at 120~m/s and 2 at
150~m/s, where the unstable pair is $3.70\pm68.2i$. Each count takes a
fraction of a second.

Figure~\ref{fig:winglocus} follows the four lowest modes from 0 to
180~m/s. The first torsion branch (93~rad/s in still air) moves down towards
the first bending frequency, turns, and crosses the imaginary axis near
70~rad/s: classical bending--torsion flutter. The flutter speed is the zero
of the spectral abscissa, found by \code{fzero}. A local Newton solution of
$\Delta(i\omega,U)=0$ in the real unknowns $(U,\omega)$ agrees to
$10^{-9}$:
\begin{equation}
  U_f=136.983977~\text{m/s},\qquad \omega_f=70.033013~\text{rad/s}.
\end{equation}
Larger windows up to 700~rad/s at $0.95\,U_f$ and $1.05\,U_f$ contain six
modes in the upper half-plane, with 0 and 1 unstable, respectively.

\begin{figure}[H]
\centering
\includegraphics[width=\textwidth]{wing_root_locus.png}
\caption{Left: modes of the continuous Goland wing from 0 to 180~m/s,
labelled by their still-air character. Right: spectral abscissa and number
of roots counted in the right half-plane.}
\label{fig:winglocus}
\end{figure}

\section{A finite-element model converges to the continuous answer}

An independent calculation follows the usual route:
\begin{itemize}
  \item cubic Hermite elements for bending and linear elements for
        torsion;
  \item projection onto the vacuum modes;
  \item air loads assembled from Theodorsen's lift and moment formulas,
        with $C(k)$ written in Hankel functions rather than the modified
        Bessel functions of the main model.
\end{itemize}
Its flutter speed decreases towards the continuous value, and the error
falls by a factor of four each time the mesh is doubled
(Table~\ref{tab:wingfe}, Figure~\ref{fig:wingconv}). This is second-order
convergence to the limit that ContourRoots computes directly.

\begin{table}[H]
\centering
\caption{Flutter speed of the finite-element model versus the continuous result.}
\label{tab:wingfe}
\begin{tabular}{rrrr}
\toprule
Elements & Modes & $U_f$ (m/s) & $|U_f^{\mathrm{FE}}/U_f-1|$\\
\midrule
8 & 8 & 137.308301 & $2.4\times10^{-3}$\\
16 & 12 & 137.065051 & $5.9\times10^{-4}$\\
32 & 20 & 137.004259 & $1.5\times10^{-4}$\\
64 & 28 & 136.989048 & $3.7\times10^{-5}$\\
128 & 36 & 136.985245 & $9.3\times10^{-6}$\\
\midrule
\multicolumn{2}{l}{continuous (ContourRoots)} & 136.983977 & ---\\
\bottomrule
\end{tabular}
\end{table}

For a tapered wing the strips do approximate the geometry, since the
properties are replaced by constant values within each strip. For a
synthetic wing whose chord decreases linearly by 30\,\%, the flutter speed
takes the values 151.58, 155.81, 156.81, 157.05 and 157.11~m/s for 2, 4,
8, 16 and 32 strips. The solution within each strip remains exact, and
there is still no modal truncation and no aerodynamic fit.

\begin{figure}[H]
\centering
\includegraphics[width=\textwidth]{wing_convergence.png}
\caption{Left: convergence of the finite-element flutter speed to the
continuous value. Right: strip refinement of a tapered wing.}
\label{fig:wingconv}
\end{figure}

\section{Time responses}

The tip transfer function, for example tip force to tip deflection, follows
from solving $K(s,U)x=f$ with the load in the tip rows. It is evaluated at
any complex $s$, so \code{cstep} and \code{clsim}
(Chapter~\ref{ch:response}) give the response directly. There is no modal
ordinary differential equation and no aerodynamic lag state.

At the high frequencies sampled by the inversion, a bending wave grows or
decays like $e^{\kappa y}$ with $\kappa=(|s|^2\mu/EI)^{1/4}$. At
$10^4$~rad/s, $\kappa L\approx27$. The transfer evaluation therefore
subdivides each strip into identical pieces with $\kappa\ell\le3$. This is
multiple shooting \citep{ascher1995}, and it is exact for constant
properties: results with $n$ and $2n$ pieces agree to $10^{-12}$.

The inversion line is placed with \code{SingularityBound=5}, to the right
of the unstable pair counted at 150~m/s. The output is expressed in mm per
kN, so that the absolute tolerance $10^{-3}$ is one micrometre.

\begin{figure}[H]
\centering
\includegraphics[width=\textwidth]{wing_time_response.png}
\caption{Left: vacuum response to a 1~kN half-sine tip pulse against an
analytic 150-mode series. Right: 1~kN tip step below and above flutter;
circles: adaptive quadrature.}
\label{fig:wingtime}
\end{figure}

Figure~\ref{fig:wingtime} shows the results. In vacuum, the pulse response
agrees with the analytic series to $1.9\times10^{-6}$~mm on a 13~mm
response. Below flutter, the tip step settles near its static deflection;
above flutter it grows at the rate $\operatorname{Re}s=3.70$ predicted by
the root count. The FFT and adaptive quadrature agree to
$7\times10^{-4}$~mm. The de Hoog method, however, converges to a plateau
0.04~mm (0.5\,\%) off while all its internal checks agree. The accelerated
series is misled by the many lightly damped poles close to the imaginary
axis. For such systems the FFT, cross-checked by quadrature, is
recommended.

\section{Scope}

The aerodynamics is two-dimensional strip theory: incompressible, with no
spanwise interaction and no stall. The structure is linear and undamped,
with no shear deformation or warping. These are modelling choices, and the
solver adds no approximation to them. Root counts are complete in the
chosen rectangles, conditional on analyticity, and time responses start
from rest. $\Delta$ is the determinant of a small dense matrix.
Matrix-valued methods, which would count characteristic values without
forming a determinant, are discussed in the development notes of the
toolbox. The same exact dynamic-stiffness idea underlies the wing analyses
of \citet{banerjee2016}.
